\scanpage{6}
\margin{$*$ Ou encore que les longueurs des dix vecteurs $a-b_i$ ($0\leq i\leq4$), $b_i-b_{i+1}$ ($0\leq i\leq4$, posant $b_5=b$) sont égales à
\(\ell=2\sin\frac{\alpha}{2}\) \textcircled{4}. Notons les longueurs correspondantes des vecteurs $a-b'_i=a+b_i$ \textcircled{$4'$} $\ell'=2\cos\frac{\alpha}{2}$}
% slanted
\margin{Notons que $\ell^2$ et $\ell'^2$ sont les solutions de l'équation quadratique $\lambda^2-4\lambda+4\sin^2\alpha=0$}

déterminé, \deleted{et l'ens} on l'appelle le \uline{sommet} de la calotte d'icosaèdre, \inserted{les autres points les \uline{points périphériques} ;} \deleted{et} Sur l'ensemble des \deleted{\unsure{directions}} \inserted{ceux-ci} (grâce à l'orientation \inserted{de $E$ donc} de $S$, \deleted{\unsure{donne|donc}} il y a alors un ordre circulaire canonique, qui en fait un torseur sous le groupe $\mathbb Z/5\mathbb Z$. Le sous-groupe \inserted{$G_C$} de $G=O(E)$ qui laisse une calotte $C$ invariante, \deleted{\illegible{1}} a dix éléments, \deleted{le sous-g} son intersection \inserted{$G'_C$} avec $G'=SO(E)$ est le groupe (can. isom. à $\mathbb Z/5\mathbb Z$) des rotations d'angle multiple de $2\pi/5=\theta$ autour du \deleted{\unsure{sommet}} \inserted{vecteur} $a$. \deleted{On obtient un} \inserted{Les cinq} points de $G_C\cap\complement G'_C$ sont \inserted{d'ordre 2, et forment avec l'élément neutre} les stabilisateurs dans $G_C$ des cinq points périphériques de la calotte ; l'élément $g_i\in G_C\cap\unsure{\complement G'_C|\complement C'_C}$ qui laisse $b_i$ invariant est la symétrie p.r. au plan engendré par $a,b_i$. La structure du groupe $G_C$ est le produit $1/2$ direct du \deleted{\unsure{\mathbb Z/5\mathbb Z}} sous-groupe $G_C$ (can. isom à $\mathbb Z/5\mathbb Z\inserted{{}=\mu_5}$) et du groupe $\mathbb Z/2\mathbb Z$, opérant sur $G_C$ par symétrie. Enfin, le groupe $G'_C\simeq\mathbb Z/5\mathbb Z=\mu_5$ opère de façon simplement transitive sur l'ens des pts périphériques de la calotte (ce qui redonne une définition de \deleted{leur} l'ordre circulaire can. sur leur ens.)

Jusqu'à maintenant le choix $\theta=2\pi/5$ n'a pas joué de rôle particulier : on aurait pu prendre $\theta=2\pi/m$, et le groupe des $m$ rotations multiples de $2\pi/m$ aurait rempli un rôle analogue. Cela va changer avec le

\uline{Lemme} L'ensemble $(x_0;x_1,u,x_{-1},x'_2,x'_{-2})$ est une calotte d'icosaèdre de sommet $u$, l'ordre circulaire canonique étant celui écrit

\uline{Démonstration} Faisant la symétrie p.r. au plan engendré par $u,x_0$ (qui échange $x_1$ et $x_{-1}$, $x'_2$ et $x'_{-2}$) on est ramené \inserted{(à une question d'orientation près)} à montrer que les \deleted{trois triangles} longueurs des trois vecteurs $x_0-x'_2,\ x_{-1}-x'_2,\ x'_2-x'_{-2}$ sont égales à la (4) -- et on on peut se dispenser de la \inserted{vérifier pour} dernier vecteur. \deleted{\unsure{De même on a}} Il sera plus astucieux de montrer que si on \deleted{\unsure{a}} \inserted{considère} l'élément \inserted{$\gamma^w_\alpha$} de $G'=SO(E)$ qui est \deleted{l'identité sur $w$, et qui dans le plan $(u,v)$ est} la rotation d'angle $\alpha$ autour de $w$ et l'élément $\gamma^u_{-\frac{\theta}{2}}$ qui est la \deleted{identité sur $u$ et qui est} \inserted{rotation} d'angle $-\theta/2$ autour de $u$, et si on pose
% archive sheet number: 5
